% RESEARCH_DRAFT_SHA256: 6a91f5df0c3e6e373524143ba3804f3b5e80f76be89c243a12a59719a86a3a01
% Layout and source-reference wrappers expanded for the web renderer.
\documentclass[11pt]{article}
\usepackage[margin=25mm]{geometry}
\usepackage{amsthm}
\usepackage{fontspec,unicode-math}
\setmainfont{Latin Modern Roman}
\setmathfont{Latin Modern Math}
\usepackage{amsmath,amssymb,mathtools,xurl,hyperref}
\hypersetup{colorlinks=true,urlcolor=blue,linkcolor=blue}
\setcounter{secnumdepth}{0}
\setlength{\parindent}{0pt}
\setlength{\parskip}{5pt}
\setlength{\emergencystretch}{3em}
\newcommand{\sref}[2]{\hyperlink{source:#1:#2}{[S#2]}}
\newcommand{\cataloguescope}[1]{\paragraph{Source record.}#1}
\newcommand{\floor}[1]{\left\lfloor#1\right\rfloor}
\newtheorem{proposition}{Proposition}
\newtheorem{lemma}[proposition]{Lemma}
\begin{document}
\paragraph{Research manuscript.} OpenAI \texttt{gpt-6-astra}, reasoning \texttt{ultra}; 11 October 2026.
\paragraph{Preserved source body.} The complete catalogue mathematical text and sources follow with layout and source links expanded, including source qualifications and its historical preparation notes. New research begins at the section headed Research attempt.

% UnsolvedMath ID 30000230; source catalogue code OWR-829-002
\section{Application-Independent Parallel Algorithms for Directed Grid Dependencies}\label{30000230}
{\small UnsolvedMath ID 30000230 · Computer Science}
\subsection{Definitions and mathematical statement}

\paragraph{Shared notation.}Standard number systems are $\mathbb N,\mathbb Z,\mathbb Q,\mathbb R,\mathbb C$. Specialized notation and index ranges follow the cited source and the definitions below.


A directed dependency graph has an arc from a computation to another computation that requires its result. A legal schedule respects these dependencies; a parallel schedule permits independent computations to proceed simultaneously.
\paragraph{Record scope and qualifications.}
The catalogue title is retained as metadata. Named certificates, contracts, bounds or frameworks in that title are not asserted as results of the original source extract.
The grid computations, permissible machine operations, correctness interface, and complexity target are not fixed.

Source-record limitation: Application-independent parallel-algorithm design is requested without a task, correctness condition, or resource target.
\paragraph{Statement recorded in the supplied source.}
Is it possible to design a uniform underlying parallel algorithm that is independent of applications but suitable for grid-based numerical computations whose data-dependency relation is nonsymmetric and can be represented accurately by a directed graph? \href{https://doi.org/10.4171/owr/2005/24}{[S2]}
\paragraph{Source status.}
The supplied catalogue labels this source record open; that metadata was checked against the supplied snapshot on 2026-10-11. The label is preserved for provenance and is not an independent verification of present open status or mathematical correctness. The catalogue literature assessment is dated 2026-08-21.
\subsection{Short English statement}
Find a reusable parallel scheduling framework for directed dependencies in numerical grids.
\subsection{Sources}
\begin{itemize}
\item[S1] Ulam AI and the UnsolvedMath Contributors, UnsolvedMath: A Curated Collection of Open Mathematics Problems, record 30000230 (OWR-829-002). Catalogue curation and metadata: CC BY 4.0. The source extract is retained; definitions, English summaries and qualifications are editorial additions.
\url{https://huggingface.co/datasets/ulamai/UnsolvedMath}.
\item[S2] Schnelle Löser für partielle Differentialgleichungen (2006), Oberwolfach Reports, DOI 10.4171/owr/2005/24, pp. 1299--1370
\url{https://doi.org/10.4171/owr/2005/24}.
\item Source-specific Oberwolfach Reports licences, including CC BY-SA 4.0 where applicable, are retained; see THIRD\_PARTY\_NOTICES.md.
\end{itemize}
\paragraph{Preparation and review.}
The original source and UnsolvedMath contributors retain attribution for the source extract. Codex (GPT-6; exact serving revision not exposed) assisted with definitions, English summaries, scope notes and formatting. No new human mathematical authorship or independently verified solution is asserted. Permanent source ID 30000230 and code OWR-829-002 checked against the complete supplied registry and current repository. Distinct upstream records, including repeated questions and broad or damaged extracts, are retained for complete coverage of the source CS category. A repeated record is not a new independent conjecture. Complete mathematical-equivalence and current-literature review has not been performed. Preparation and final review follow contribution prompts C and D, with source gaps disclosed. Structural and rendering checks do not establish a new theorem or open status.
\label{end:30000230}

\section{Research attempt}
\subsection{Literature, source context and scope}
The full original catalogue body and source qualifications are preserved above. The relevant primary workshop passages were read in the cached publisher report. Historical claims and subsequent attributed results are distinguished from the mathematical work below. A restricted example, standard lemma or known theorem is not counted as new progress on the entire original target. Missing assumptions and repeated or rhetorical extracts are not silently turned into different solved problems.

\subsection{Testing the dependency-graph route on unit tasks}
The report already proposes directed dependency graphs, ready-task sweeps and priorities. We make that existing route precise for a restricted class: a finite directed acyclic graph of $W$ unit-duration tasks, with each edge requiring its tail to finish before its head starts. There are $p\ge1$ identical processors, and $L$ is the number of vertices on a longest directed path. At each integer time, execute any $p$ ready tasks, or all ready tasks if fewer than $p$ remain ready.
\begin{proposition}
This list schedule terminates after $T\le W/p+(1-1/p)L$ time steps. Every feasible schedule requires at least $\max(W/p,L)$ steps.
\end{proposition}
\begin{proof}
A nonempty finite acyclic graph has a source: otherwise repeatedly following an incoming edge would revisit a vertex and form a cycle. Consequently the schedule always executes a task until it terminates. Call a step full if it executes $p$ tasks and underfull otherwise, and write their numbers as $F,U$. Every underfull step removes all sources of the remaining graph. Removing all sources reduces its longest-path length by at least one: the first vertex of any path remaining afterward had an incoming ancestor before the removal, and a chain to a source prepends at least one vertex to that path. Thus $U\le L$. Every step executes at least one task, so $W\ge pF+U$. Therefore
\[
T=F+U\le W/p+(1-1/p)U\le W/p+(1-1/p)L.
\]
The workload bound follows because at most $p$ tasks finish per step; the path bound follows because all $L$ tasks on a longest path must be executed in order.
\end{proof}
For the $n$ by $n$ grid with each task depending on its existing north and west neighbors, $W=n^2$ and $L=2n-1$. The path length follows because each edge increases the sum of the two coordinates by one, and a corner-to-corner monotone path attains that length. This gives a concrete performance certificate once dependencies and costs are supplied.
\subsection{Verification and remaining gaps}
A deterministic list schedule on the $3$ by $3$ grid with two processors was computed and compared with both bounds. The attempted extension to arbitrary application graphs stops at cycles: the two tasks $a\to b\to a$ have no ready task, although an iterative numerical interpretation might still be meaningful. Such an interpretation requires convergence and update semantics absent from a precedence graph. Heterogeneous runtimes, communication, dynamic dependencies and application-specific cyclic iteration are not handled. This standard work/path estimate is not a new general framework or completion of the report's broad agenda.

\section{References and source attribution}
Primary workshop report: \url{https://doi.org/10.4171/owr/2005/24}. Relevant passages read: 2005/24 pp. 1359--1360, dependency graphs, sweeping, priority and cyclic applications.

All retrievals were checked on 11 October 2026. Original question attribution remains in the preserved Sources section. External theorems are attributed and not claimed as new discoveries or fully verified merely by citation.


\paragraph{Integration review.} Codex (GPT-6; exact serving revision not exposed) reviewed the written arguments and expanded layout and source-reference wrappers. Repository metadata and the production web renderer were checked using Python 3.13 and JavaScript with MathJax 3.2.2. This model-assisted review is not independent human proof verification; no PDF compilation or formal proof checking is asserted.

\section{Completion Estimate}
\noindent\textbf{COMPLETION: 0\%}
No substantiated new progress toward the entire original target is claimed. The worked results and failed route are delimited above; a known result, illustrative model, or auxiliary obstruction is not counted as original-target completion.

\section{Final status and attribution}
\textbf{Attempt status: unresolved.} This describes this research attempt, not a new ruling on the source catalogue's status or on external literature claims.
The mathematics and this self-audit were generated by OpenAI \texttt{gpt-6-astra}, reasoning \texttt{ultra}, on 11 October 2026. Codex (GPT-6; exact serving revision not exposed) orchestrated the batch and contributed formatting/repository review. No human mathematical authorship, independent human proof checking, or formal verification is claimed. The historical preparation paragraphs in the reproduced catalogue body concern that body only.

Contribution prompt A was applied to the complete definition and supplied source record, followed by the adversarial checks of prompt B. Local reproduction files are \texttt{batch016/checks.py}, \texttt{computations.json}, and the per-ID audit JSON. They were not public links at preparation time. Numerical checks are finite and do not prove unrestricted results. No LaTeX engine was found locally; compilation was not performed by this research agent.

\end{document}
